\begin{table*}[t]
\caption{Related-work positioning. BodyShield is framed as falsification-guided embodiment-control repair, not as a replacement for broad randomized training, formal safety control, or cross-embodiment foundation policies.}
\label{tab:related-work-positioning}
\centering
\footnotesize
\begin{tabular}{p{0.18\linewidth}p{0.36\linewidth}p{0.36\linewidth}}
\toprule
Area & Prior focus & BodyShield focus \\
\midrule
Domain randomization & Train over sampled randomized simulator parameters. & Search for low-cost perturbations that actually break the current policy, then repair against discovered counterexamples. \\
Robust/adversarial RL & Train against disturbances or adversarial forces. & Treat latency, calibration, actuation limits, sensing shifts, payload, tool geometry, and contact as auditable embodiment-control axes. \\
Counterexample-guided repair & Repair safety or state-space violations found by a verifier. & Repair hidden body/control assumptions and require oracle feasibility for claimed breaking perturbations. \\
SysID and retuning & Estimate model parameters and retune control. & Diagnose which body/control changes expose brittleness even when the task remains feasible. \\
CBF/MPC/reachability & Provide controller-level safety or robust feasibility guarantees. & Provide empirical falsification and repair machinery; no formal safety certificate is claimed. \\
Cross-embodiment policies & Train broad policies across robots or datasets. & Supply a stress-test and repair protocol for scoped embodiment shifts, without claiming foundation-policy generality. \\
\bottomrule
\end{tabular}
\end{table*}
